\documentclass[10pt,a4paper]{amsart}
\usepackage[margin=19mm]{geometry}
\usepackage{amsmath,amssymb,mathtools}
\usepackage[scr=rsfs,cal=euler]{mathalfa}
\usepackage{xfrac}
\usepackage[unicode,hidelinks,bookmarks=false]{hyperref}
\newcommand{\ie}{i.e.\ }
\newcommand{\cf}{cf.\ }
\newcommand{\resp}{respectively\ }
\newcommand{\Subsection}{\subsection}
\providecommand{\nobreakdash}{\nobreak\hbox{-}}
\setlength{\emergencystretch}{2em}
\multlinegap0pt
\usepackage{bm}
\usepackage{bbm}
\usepackage{dsfont}
\usepackage{xspace}
\usepackage{cancel}
\usepackage{mathtools}
\usepackage[shortlabels]{enumitem}
\usepackage{amsthm}
\makeatletter
\@ifundefined{theorem}{\newtheorem{theorem}{Theorem}}{}
\@ifundefined{lemma}{\newtheorem{lemma}[theorem]{Lemma}}{}
\makeatother
\makeatletter
\newcommand{\Mat}{\operatorname{Mat}}
\newcommand{\diag}{\operatorname{diag}}
\expandafter\ifx\csname fthm\endcsname\relax
\newenvironment{fthm}
  {\begin{mdframed}[innertopmargin = 3pt, innerbottommargin=7pt,skipabove=5pt,skipbelow=5pt,linewidth=0.25pt,nobreak=true,align=center]\begin{ftheo}}
\fi
\makeatother
\title[Vanishing of higher Legendrian homology for rainbow closures]{Vanishing of higher Legendrian homology for rainbow closures}
\author{Roger Casals, Alexander Simons}
\date{Source: arXiv:2608.24844v1 (CC BY 4.0)}
\begin{document}
\maketitle

\noindent\emph{Source and licence.} Vanishing of higher Legendrian homology for rainbow closures. Version arXiv:2608.24844v1 (CC BY 4.0), distributed under the Creative Commons Attribution 4.0 International licence, as recorded in the arXiv licence metadata for this exact version and shown on the abstract page https://arxiv.org/abs/2608.24844v1. Full rights notice: licenses/arxiv-2608.24844-CC-BY-4.0.txt.\par\smallskip

\noindent\textit{Editorial scope.} Counted unit: the Lemma whose source label is \texttt{lem:LU\_decomposition} in the article (source0.tex lines 386--395, source label \texttt{lem:LU\_decomposition}), delivered together with the article's own complete original proof of it (source0.tex lines 397--471, one \texttt{proof} environment of 75 lines). No mathematics is written, repaired or re-derived: statement and proof are the article's own text line for line. Required context retained uncounted: none. Every definition, display and label of those lines is retained unchanged. Omitted from this package and not redistributed: 1--385, 396--396, 472--1017. Nothing inside a retained excerpt is omitted or altered: every retained body is delivered exactly as the article's lines are written, so no omission receipt of any kind accompanies this package. Citations: the retained excerpts carry no citation command at all, so no bibliography entry of the article is transcribed and no reference is redistributed. Reference adaptation: none was needed and none was made; every reference inside the delivered unit resolves inside it, so no unresolved reference or citation is produced. Printed slips kept literal: no printed word, symbol, hypothesis or number of the counted statement or of its proof was altered. Layout and class: the article's own class is \texttt{article} with packages including fontenc, inputenc, lmodern, microtype, geometry, amsmath, amssymb, amsthm, mathtools, mathrsfs, array, booktabs, enumitem, xcolor; the wrapper is the lane's pinned \texttt{amsart} editorial wrapper at 10pt on A4, which restates the theorem-like environments the retained text needs and adds the article's own macro definitions that the retained text needs (\texttt{\textbackslash Mat}, \texttt{\textbackslash diag}), with extra wrapper packages (bm, bbm, dsfont, xspace, cancel, mathtools, enumitem). The delivered numbering is the unit's own. Assets: no retained span contains a figure environment, a graphics-inclusion command, a PDF-page inclusion, a table or a TikZ picture; no image file is copied into this package, the package declares no asset dependency and no image file is redistributed with it. Licence: version arXiv:2608.24844v1 of this article is distributed under the Creative Commons Attribution 4.0 International licence, as recorded in the arXiv licence metadata for this exact version and shown on the abstract page https://arxiv.org/abs/2608.24844v1. A full rights notice is delivered at licenses/arxiv-2608.24844-CC-BY-4.0.txt. Characters: every retained excerpt is pure ASCII, so the delivered engine, XeLaTeX with its default Latin Modern text font, reports no missing character.
\begin{lemma}\label{lem:LU_decomposition}
Let \(S\) be a unital, not necessarily commutative, ring, and let
\(Y\in\Mat_n(S)\). Then the following are equivalent:
\begin{enumerate}[label=\textup{(\roman*)},leftmargin=2.2em]
\item Every leading principal submatrix of $Y$ is invertible.
\item There is a factorization $Y=D_YL_YU_Y$, where \(D_Y\) is invertible diagonal, \(L_Y\) is lower
unitriangular, and \(U_Y\) is upper unitriangular.
\end{enumerate}
In addition, if such factorization exists, it must be unique.
\end{lemma}

\begin{proof}
To ease notation we denote by $ Y^{[r]}$ the $r\times r$ leading principal submatrix of a matrix $Y$. Let us first show $(ii)\Longrightarrow(i)$, so we assume that the factorization \(Y=D_YL_YU_Y\) exists.  Since the factors are respectively
diagonal, lower triangular, and upper triangular, taking a leading
principal block commutes with their product. There we obtain
\[
 Y^{[r]}=D_Y^{[r]}L_Y^{[r]}U_Y^{[r]}.
\]
Every factor on the right is invertible, so \(Y^{[r]}\) is invertible, and $(ii)\Longrightarrow(i)$ is proven.\\

For $(i)\Longrightarrow(ii)$, we suppose that all the leading principal submatrices \(Y^{[r]}\) are invertible.  We argue by
induction on the size \(n\in\mathbb{N}\).  The case \(n=1\) holds tautologically. Let us write
\[
 Y=\begin{pmatrix}A&b\\ c&x\end{pmatrix},
\]
where \(A=Y^{[n-1]}\).  By the induction hypothesis, there is a unique factorization
\[
 A=D'L'U'.
\]
By setting $s:=x-cA^{-1}b$, note that the matrices
\[
 P:=\begin{pmatrix}I&0\\-cA^{-1}&1\end{pmatrix},
 \qquad
 Q:=\begin{pmatrix}I&-A^{-1}b\\0&1\end{pmatrix}
\]
are invertible and satisfy
\[
 PYQ=\begin{pmatrix}A&0\\0&s\end{pmatrix}.
\]
Since \(P,Y,Q\) are invertible, the block-diagonal matrix
\(\diag(A,s)\) is invertible.  Multiplying by
\(\diag(A^{-1},1)\) shows that \(\diag(\mbox{Id},s)\) is invertible. The
bottom-right equations for a two-sided inverse then give a two-sided
inverse for \(s\), and thus \(s\) is invertible.
Set
\[
 d_n:=s,
 \qquad
 u:=(D'L')^{-1}b,
 \qquad
 \ell:=s^{-1}c(U')^{-1},
\]
and define the following extensions of $D',L'$ and $U'$:
\[
 D_Y:=\begin{pmatrix}D'&0\\0&s\end{pmatrix},
 \quad
 L_Y:=\begin{pmatrix}L'&0\\\ell&1\end{pmatrix},
 \quad
 U_Y:=\begin{pmatrix}U'&u\\0&1\end{pmatrix}.
\]
Then the upper-left, upper-right, and lower-left blocks of
\(D_YL_YU_Y\) are respectively \(A,b,c\).  Its lower-right entry is
\begin{align*}
 s(1+\ell u)
 &=s+c(U')^{-1}(D'L')^{-1}b\\
 &=s+cA^{-1}b=x,
\end{align*}
where
\((U')^{-1}(D'L')^{-1}=A^{-1}\).  Thus \(Y=D_YL_YU_Y\), which proves  $(i)\Longrightarrow(ii)$.\\

\noindent For uniqueness, we suppose that \(Y=D_YL_YU_Y\) and proceed by induction, as above.  The induction hypothesis is that leading
\((n-1)\times(n-1)\) block determines \(D',L',U'\).  The
upper-right block then determines
\[
 u=(D'L')^{-1}b.
\]
Using the lower-left and lower-right submatrices we obtain
\[
 x=d_n+c(U')^{-1}u,
\]
so
\[
 d_n=x-c(U')^{-1}u=x-cA^{-1}b=s.
\]
and we also have $\ell=d_n^{-1}c(U')^{-1}$. Thus every factor is uniquely determined.
\end{proof}

\end{document}
